% =====================================================================
%  chunks/rev_ch03.tex —— 第 3 章 读书笔记（tr-ji 撰）
% =====================================================================
\chapterreview{%
第 3 章占全书近三分之一的篇幅，是全书的代数拓扑主体。章标题本身就是一份纲领：单纯复形与 $CW$ 复形，同调与上同调，它们与同伦论的关系，障碍。四句话对应四层任务，也使本章自然分成三段：前五节造机器，第六、七节用机器算同伦，第八、九节换一台更好的机器。第 4 章的分类定理之所以能够陈述，靠的全是这里造出的工具。

第一段的重心是把「数洞」变成代数。单纯复形给出链复形 (2.2) 与同调群 (2.3)，上同调的乘法由对角映射拉回定义 (2.20)，再靠卡积把上同调送回同调 (2.30)；到了闭可定向流形，基本类的卡积恰好就是 Poincaré 对偶 (2.31)。此后三步都是松绑：相对同调与 $CW$ 复形把计算归约到胞腔，奇异同调让任意空间都能算，层与局部系数让系数真正依赖底空间。到第五节，系数已是模，Reidemeister 挠 (5.14) 与简单同伦型登场，「同伦等价能否由初等操作实现」第一次成为与同调并列的问题。

第二段是组装机器。§6 用障碍上链回答延拓问题：纤维的 $\pi_q$ 被组织成 $H^{q+1}(B;\pi_q(F))$ 中的一个类，截面存在当且仅当这个类为零，分类问题于是变成上同调计算。再由万有丛得到分类定理：$B$ 上丛的等价类就是 $[B,B_G]$（定理 6.1）。同一节把 Eilenberg--MacLane 表示成基本类的拉回 (6.6)，把「自然变换」具体化为 Steenrod 代数；Adams 谱序列随即把稳定同伦群的计算交给 $\Ext_{\mathcal{A}_p}^{\ast,\ast}$ (6.16)。§7 把这条流水线补全：Leray 谱序列用纤维丛的局部系数数据逐层逼近总空间 (7.2)，Postnikov 塔则把一个空间拆成一串 $k$-不变量 (7.11)。

第三段视角最高。$K$ 理论与配边论说明同调论并不唯一：Bott 周期性 (8.6) 使复情形以 $2$、实情形以 $8$ 为周期，Thom 同构 (8.21) 与 Riemann--Roch 型定理把示性类接到指标；复上配边论更进一步，$\Omega_U^*$ 自身带一个形式群 (9.10)，其对数正是 $g(u)=\sum_n\CP^n u^{n+1}/(n+1)$，于是 $\Omega_U^*(\CP^{\infty})$ 的信息全压进一个级数 (9.12)。

两处最易失足。Reidemeister 挠只在 $\rho$-零调复形上有定义（定义 5.1），且取值模 $\pm\zeta^{\alpha}$，把它当作一般复形的不变量是常见误读。另一处是谱序列的读法：$E_2$ 只是第一近似，微分与扩张必须逐案处理；(7.7) 那种「$E_2\cong H^*(X\times F)$」只在 $\pi_1$ 平凡作用时成立。

若只带走一句话：本章的主线是\emph{同调不唯一}。前七节把经典同调论做到极致，第八、九节却换一套广义同调论，算出前者看不见的东西。球面稳定同伦群里那些具体关系，例如 $240\sigma=0$ 与 $\eta^3=12\nu$ (7.15)，只是序曲；真正的收获是形式群。Adams--Novikov 谱序列 (9.9) 把 Adams 谱序列中的 $H\Z/p$ 换成 $MU$，这一替换正是形式群论进入稳定同伦论的入口，也是 Novikov 本人的主战场。作者把圆周作用与 Atiyah--Bott 型公式 (9.20) 放在全章末尾，说明这一节不是附录，而是全书研究前沿的起点。
}

% CHUNK rev_ch03 DONE
